Toss a fair coin forever. The proportion of heads tends to one half, but how fast, how rarely does it stray far, and what
shape does the whole path take when you look from far away? This book answers with the great limit theorems of probability
in their sharp forms: Donsker's theorem, which turns every random walk into Brownian motion; Berry and Esseen's bound on the
error of the central limit theorem; the stable laws that replace the bell curve when variances are infinite; Stein's method
of approximation; Cram\'er's large deviations; the concentration of measure that makes high dimensions predictable; de
Finetti's theorem on exchangeable sequences; and Birkhoff's ergodic theorem, which explains the digits of continued fractions
and proves that dense sets of integers contain progressions. It ends with the Poisson--Dirichlet law, which governs both the
cycles of a random permutation and the prime factors of a random integer. Every result is proved in full, with one clearly marked exception from analytic number theory. Each chapter
opens with the story behind its ideas, ends with exercises and complete solutions, and comes with a Python notebook.
